\chapter{Introduction}
\label{chap:introduction}

YouTube is one of the largest public spaces for news, opinion and debate: it has more than 2,500 million monthly users \autocite{statista2025social}, and about a third of adults in the United States regularly get news from it \autocite{pewresearch2025social}. For researchers who study public opinion, misinformation or polarization, such as the PROMUEVA project at Universidad del Valle \autocite{promueva_project}, the videos and comments published on the platform are a valuable source of evidence. Collecting that evidence in a way that other researchers can examine, however, is harder than it looks.

The tools available today fall into two groups. Commercial social listening platforms collect and analyze data at scale, but their methods are proprietary, they do not expose the raw data or the processing behind their figures, and their licenses are out of reach for most academic projects. Open-source and academic tools are free and transparent, but they leave researchers to translate their questions into the syntax of the YouTube Data API by hand and to document each collection themselves. On top of that, the platform itself limits what can be reproduced: search results for the same query change over time \autocite{efstratiou2025youtube}, and the YouTube API Services Developer Policies allow API data to be stored for at most 30 days \autocite{youtube_api_services_developer_policies}. A dataset collected from the API cannot be obtained again identically.

This thesis presents OpenTube Insights, a web application that automates the collection, processing, analysis and visualization of YouTube data for academic research, and that treats the collection procedure, rather than the dataset, as the unit another researcher can inspect and repeat. A researcher describes a search in natural language; a pre-trained large language model drafts the API parameters, which the application validates and the researcher reviews before saving them. Every collection run records its exact parameters and every request it makes, including its quota cost, so the collection can be audited after the data itself has been deleted. The collected comments are normalized by versioned rules and labeled by a sentiment model whose name and version are stored with every result, and the results are presented as summaries, charts, filters and exports that state the period and the model behind each figure.

The application was built with Django REST Framework, PostgreSQL and React, packaged with Docker Compose, and developed in two-week sprints following an adapted Scrum framework, with an AI coding agent implementing each user story under a documented plan, review and verification cycle. Its tests cover each component without calling external services, and the complete workflow was validated on two real topics from the context of the PROMUEVA project: the MIO mass transit system and the 2021 National Strike in Cali.

The rest of the document is organized as follows. Chapter~\ref{chap:problem_statement} states the problem and Chapter~\ref{chap:justification} justifies it from an academic, economic and social point of view. Chapter~\ref{chap:objectives} sets the general and specific objectives, and Chapter~\ref{chap:project_scope} the scope of the project. Chapter~\ref{chap:reference_framework} presents the reference framework: the key concepts, the state of the art and the theoretical framework, including the ethical and legal constraints of working with YouTube data. Chapter~\ref{chap:methodology} describes the methodology, and Chapter~\ref{chap:development} reports the development of the application, one section per planned activity, from the design of the architecture to its testing and validation. Chapter~\ref{chap:conclusions} closes with the conclusions, the limitations and future work, and Appendix~\ref{chap:annex_screens} collects screenshots of the application.
